\chapter{May}

\section[May 4th]{May 4\textsuperscript{th}}
The class started today with some visual puns, as most of the class time was spent on a speech about child protection by a police officer. It was really interesting, especially regarding the role of internet monitoring in preventing the abuse of children by strangers. Furthermore, the examples narrated by the officer were surprising; it is weird to think that such things can happen so close to where I live.

An important point made to the future teachers in the audience was that sometimes the danger comes from the parents; in those cases, the trusted adults would be the teachers themselves.

Unfortunately, they didn't follow the scheduled time. It was supposed to be from 7:00 pm to 8:00 pm, but it actually went until about 9:00 pm. Because of this, we couldn't discuss much about Academic Writing today.

Either way, we were able to review some more conjunctions of agreement, example, support, and emphasis. It is always nice to improve our repertoire of words and ways to express our ideas. It shows that many times such connectives cannot be translated to Portuguese easily; we end up only accessing them once we are able to think in English while writing.

Some articles that I found:
\begin{itemize}
    \item \url{https://arxiv.org/pdf/2602.12891}
    \item \url{https://arxiv.org/pdf/2205.12615}
    \item \url{https://muse.jhu.edu/pub/1/article/971097/pdf}
    \item \url{https://plato.stanford.edu/entries/computer-science/}
    \item \url{http://projecteuclid.org/journalArticle/Download?urlid=10.1305%2Fndjfl%2F1093893717}
    \item \url{https://periodicos.sbu.unicamp.br/ojs/index.php/manuscrito/article/view/8668955/28277} (My favorite for now)
\end{itemize}

\section[May 11th]{May 11\textsuperscript{th}}

We had a brief encyclopedia reading session from the book ``There's no Such Thing as a Silly Question''. We started by reading a question about the existence of aliens, learning about the size of the universe and other galaxies. The second question was whether animals can commit crimes (it depends on the place and era, although currently, in most jurisdictions, animals cannot be convicted of crimes), and the last was about how small a star can be (about the size of the planet Saturn).

I like this method better than the other children's books; it feels more useful and dynamic than the stories.

After that, we had an activity using physical dictionaries, where we searched for transition words and phrases and created examples for each transition. It was nice, albeit a tad time-consuming. I felt that the efficiency of learning dropped slightly. It took a while to search for each word, and since only one member of the team had a pen, we needed to write one at a time. On a positive note, it was great for teamwork and getting acquainted with my colleagues; if that was the intention, then I would say it was successful.

\section[May 18th]{May 18\textsuperscript{th}}

There was no class today, as it was the ``Semana Acadêmica de Letras''.

The title of the speech was ``SAÚDE CORPORAL E MENTAL \dots\ E A VOZ? BREVE HISTÓRIA DO CORPO, DAS EMOÇÕES E DA VOZ''. The content presented by Dr. Dobrić was really interesting, especially regarding the health and care we need to take with our voices. If we neglect it, bad consequences can arise. This is even more important for my colleagues in this class who are teachers and use their voices to work with their students.

\section[May 25th]{May 25\textsuperscript{th}}

The book of the day was ``If You Want to See a Whale'', a book about focusing on your objectives, not getting distracted, and not giving up so that you will eventually find what you want. This time, it took about 25 minutes to finish the story session (6:45 pm to 7:10 pm). There was some reflection about how it related to Academic Writing, specifically regarding patience, persistence, determination, and so on. I can see it, but the correlation seems a bit vague; we can relate it to almost anything in life that demands patience.

Then we learned about the ``hedging device'', in Portuguese \textit{dispositivo modalizador}. It is a way of not committing fully to an affirmation, for example, using clauses like ``it appears'', ``the narrative suggests'', and so on.

One positive point is that the book served as a way to work with both creative and academic writing. ``If you want to see a whale, you will need time'' becomes ``The narrative suggests that meaningful observation requires temporal openness.'' This helps provide a more literal way of expressing the ideas of the text. The first phrase sounds simpler while being highly metaphorical. The latter is less interpretative; the idea is clearly and objectively presented alongside who stated it (the narrative).

Then the professor talked about voice. In forensic linguistics, we can identify the ``text voice''; she explained how we can identify someone by their style of writing. We also discussed the physical voice and how we can keep our speech healthy.

After all this, we were asked again: what is academic writing? Of course, it is the style of writing used in academic publications. As a form of writing, it presents certain characteristics: it is impersonal, objective, precise, concise, formal, structured, focused, and backed up by evidence. The primary goal is to properly communicate what was done during the research and its results. Furthermore, communicating the research and results shows how it is actually useful and how it contributes to the field in which the writing is being published.

Some nice tips that are important for me are:
\begin{itemize}
    \item Do not use contractions.
    \item Use active voice.
    \item Avoid phrasal verbs.
    \item Use the most appropriate formal negative forms (instead of ``didn't yield any results'', use ``yielded no results'').
    \item Do not use ``etc.'' or ``and so forth''; it shows that you don't really know what you are talking about. Be more specific, using phrases like ``and other electronic devices'' or ``such as'' followed by examples.
\end{itemize}

When creating an abstract, we can follow 5 moves. The interesting part is that the abstract follows the structure of the article itself in the same order, albeit in a very summarized manner using just a few sentences.
